\documentclass[conference]{IEEEtran}
\IEEEoverridecommandlockouts
% The preceding line is only needed to identify funding in the first footnote. If that is unneeded, please comment it out.
\usepackage{cite}
\usepackage{amsmath,amssymb,amsfonts}
\usepackage{algorithmic}
\usepackage{graphicx}
\usepackage{textcomp}
\usepackage{xcolor}
\usepackage{booktabs}
\usepackage{subcaption}
\def\BibTeX{{\rm B\kern-.05em{\sc i\kern-.025em b}\kern-.08em
    T\kern-.1667em\lower.7ex\hbox{E}\kern-.125emX}}

\begin{document}

\title{Deep Generative Models for Density Estimation and Unpaired Image Translation: A Comprehensive Study of MAF and CycleGAN\\
{\footnotesize \textsuperscript{*}Note: Assignment 2 for Deep Generative Models Course}
}

\author{\IEEEauthorblockN{Mohammad Taha Majlesi}
\IEEEauthorblockA{\textit{Department of Computer Engineering} \\
\textit{University of Tehran}\\
Email: taha.majlesi@ut.ac.ir}
}

\maketitle

\begin{abstract}
This paper investigates two complementary deep generative paradigms: (i) Masked Autoregressive Flow (MAF), a likelihood-based normalizing flow enabling explicit density estimation and anomaly scoring, and (ii) CycleGAN, an adversarial framework for unpaired image-to-image translation leveraging cycle consistency and identity constraints. We present unified theoretical derivations including change of variables formulations, Jacobian determinant computations, and autoregressive factorizations. Architectural design rationales are detailed for MADE masking schemes, ResNet-based generators, and PatchGAN discriminators. Training methodologies encompass optimization objectives and replay buffer stabilization techniques. Empirical validation is conducted through density modeling on industrial capsule imagery (MVTec AD dataset) and stylistic horse-zebra translation, evaluated via AUROC metrics (0.72 achieved) for anomaly detection and qualitative assessments for translation fidelity. Critical analysis addresses fundamental limitations including likelihood calibration challenges and geometric transformation constraints. Ethical deployment considerations and future research directions are systematically examined. This work consolidates foundational theory, implementation specifics, experimental results, and actionable insights for practitioners seeking robust generative modeling solutions.
\end{abstract}

\begin{IEEEkeywords}
Deep generative models, normalizing flows, masked autoregressive flow, CycleGAN, unpaired image translation, anomaly detection, cycle consistency, density estimation
\end{IEEEkeywords}

\section{Introduction}
\IEEEPARstart{D}{eep} generative models have revolutionized the ability to learn complex data distributions and perform cross-domain transformations. Flow-based models, exemplified by Masked Autoregressive Flows (MAF), provide tractable likelihood computation essential for statistical inference tasks including anomaly detection and density estimation. Adversarial translation frameworks, particularly CycleGAN, emphasize perceptual realism in conditional generation without requiring paired supervision. This paper examines MAF and CycleGAN as representative instances of likelihood-based and adversarial paradigms respectively, providing technically rigorous yet implementation-focused analysis suitable for reproducible experimentation.

\subsection{Motivation and Context}
Industrial defect detection systems benefit from explicit probability density estimation, enabling prioritization of data regions exhibiting atypical statistical signatures. Manufacturing quality control demands automated anomaly identification with quantifiable uncertainty. In artistic, biological, and medical imaging domains, scarcity or infeasibility of paired training examples necessitates unpaired translation frameworks. Cross-modal synthesis (e.g., MRI-CT translation, seasonal style transfer) requires learning mappings between domains without pixel-level correspondence. Understanding both likelihood-based and adversarial approaches enables hybrid architectures combining their complementary strengths—explicit probabilistic scoring with perceptually realistic generation.

\subsection{Paper Contributions}
This work provides comprehensive treatment of two foundational generative modeling paradigms through:
\begin{itemize}
 \item Unified mathematical derivations for coupling and autoregressive transformations with computational complexity analysis and practical implementation implications
 \item Detailed architectural specifications including layer-by-layer parameter counts, activation functions, normalization schemes, and design trade-offs
 \item Structured evaluation methodology employing threshold-independent metrics (AUROC) for anomaly detection and qualitative/quantitative assessments for image translation
 \item Systematic dissection of CycleGAN training objectives, stability mechanisms (replay buffers, learning rate scheduling), and inherent constraints on geometric transformations
 \item Critical examination of failure modes, likelihood calibration challenges, and ethical deployment considerations
 \item Reproducibility documentation with complete hyperparameter specifications, preprocessing pipelines, and computational resource requirements
\end{itemize}

\section{Background}
\subsection{Normalizing Flows: Mathematical Foundation}
Normalizing flows learn an invertible mapping $f_{\theta}: X \to Z$ pushing data $x$ to latent $z$ under a simple base density $p_Z$ (typically isotropic Gaussian). The log-likelihood decomposes across successive bijections:
\begin{equation}
\log p_{\theta}(x) = \log p_Z(f_{\theta}(x)) + \log \left| \det \frac{\partial f_{\theta}(x)}{\partial x} \right|.
\end{equation}

\paragraph{Composing Flows} Given sequence $f = f_K \circ \cdots \circ f_1$, Jacobians multiply yielding:
\begin{equation}
\log p_{\theta}(x) = \log p_Z(z_K) + \sum_{k=1}^{K} \log \left| \det \frac{\partial f_k}{\partial z_{k-1}} \right|, \quad z_k = f_k(z_{k-1}).
\end{equation}
Each bijection $f_k$ must have tractable Jacobian determinant computation (ideally $\mathcal{O}(D)$ rather than $\mathcal{O}(D^3)$ for full matrix determinant).

\paragraph{Invertibility Requirements} Training via maximum likelihood requires forward pass (data$\to$latent) for density evaluation; sampling requires inverse (latent$\to$data). Coupling layers provide both directions efficiently; autoregressive flows favor one direction.

\subsection{Autoregressive Modeling and MADE}
Autoregressive factorization decomposes joint density:
\begin{equation}
p(x) = \prod_{d=1}^{D} p(x_d \mid x_{<d}),
\end{equation}
where each conditional $p(x_d \mid x_{1:d-1})$ may be parameterized by neural networks. MADE (Masked Autoencoder for Distribution Estimation) enforces this dependency structure via binary masks $M$ applied to weight matrices, zeroing connections from $x_j$ to $x_i$ whenever $j \geq i$ under the chosen ordering.

\paragraph{Masking Strategy} Assign each hidden unit $h$ a random number $m(h) \in \{1, \dots, D-1\}$ representing maximum input dependency. Mask from input $x_d$ to hidden $h$ is 1 iff $d \leq m(h)$; mask from hidden $h$ to output parameter for $x_d$ is 1 iff $m(h) < d$. This guarantees autoregressive property.

\paragraph{Connection to MAF} MAF stacks MADE blocks where each outputs transformation parameters $(\mu_d, \sigma_d)$ depending autoregressively on previous dimensions. Permuting dimension order between blocks increases expressiveness by varying dependency patterns.

\subsection{Adversarial Image Translation}
Generative adversarial networks pit generator $G_{\theta}$ against discriminator $D_{\phi}$ in a minimax game:
\begin{equation}
\min_{\theta} \max_{\phi} \mathbb{E}_{x \sim p_{\text{data}}}[\log D_{\phi}(x)] + \mathbb{E}_{z \sim p_Z}[\log(1 - D_{\phi}(G_{\theta}(z)))].
\end{equation}

\paragraph{Conditional GANs (cGANs)} Conditioned on input $x$, generator produces $G(x)$ targeting output domain; discriminator evaluates pairs $(x, y)$. Pix2pix exemplifies supervised paired translation with $\mathcal{L}_1$ reconstruction penalty.

\paragraph{Unpaired Translation Challenge} CycleGAN addresses absence of paired $(x,y)$ by learning dual mappings $(G: X \to Y, F: Y \to X)$ with cycle consistency constraint $F(G(x)) \approx x$ and $G(F(y)) \approx y$, preventing mode collapse while avoiding explicit correspondence.

\subsection{Instance Normalization vs Batch Normalization}
Instance normalization (IN) normalizes each sample independently across spatial dimensions:
\begin{equation}
\text{IN}(x) = \gamma \frac{x - \mu(x)}{\sigma(x)} + \beta,
\end{equation}
where $\mu, \sigma$ computed per-instance per-channel. This contrasts with batch normalization (BN) using mini-batch statistics. IN better preserves individual image style statistics, critical for style transfer tasks.

\subsection{Residual Learning}
Residual blocks learn additive corrections $\mathcal{F}(x)$ to input:
\begin{equation}
y = x + \mathcal{F}(x),
\end{equation}
mitigating vanishing gradients in deep networks. For image translation, residuals encode localized texture modifications while identity path maintains structure.

\section{Methodology: Detailed Implementation}
\subsection{MAF Architecture Specification}
\subsubsection{MADE Block Structure}
Each MADE block consists of:
\begin{itemize}
\item \textbf{Input layer}: $D = 49{,}152$ dimensions (flattened $128 \times 128 \times 3$ RGB images).
\item \textbf{Hidden layers}: Two fully-connected layers with 512 units each, masked weights enforcing autoregressive property.
\item \textbf{Output layer}: $2D$ outputs encoding per-dimension mean $\mu_d(x_{<d})$ and log-scale $\log \sigma_d(x_{<d})$.
\item \textbf{Activation}: ReLU for hidden layers; linear output.
\end{itemize}

\subsubsection{Transformation Equation}
For dimension $d$ in block $k$:
\begin{equation}
z^{(k)}_d = \frac{x^{(k)}_d - \mu^{(k)}_d(x^{(k)}_{<d})}{\sigma^{(k)}_d(x^{(k)}_{<d})},
\end{equation}
where $\mu^{(k)}_d, \sigma^{(k)}_d$ are neural network outputs. Inverse for sampling:
\begin{equation}
x^{(k)}_d = \mu^{(k)}_d(x^{(k)}_{<d}) + \sigma^{(k)}_d(x^{(k)}_{<d}) \cdot z^{(k)}_d.
\end{equation}

\subsubsection{Multi-Block Composition}
Seven MAF blocks stacked with dimension permutations (reversal or random shuffle) between blocks. Final latent $z^{(7)}$ compared against standard Gaussian base. Total trainable parameters: approximately 18M (dominated by $512 \times 49{,}152$ weight matrices).

\subsubsection{Training Objective and Gradients}
Minimize negative log-likelihood:
\begin{equation}
\mathcal{L}(\theta) = -\frac{1}{N} \sum_{i=1}^{N} \left[ \log p_Z(z_i^{(K)}) + \sum_{k=1}^{K} \sum_{d=1}^{D} \log \sigma^{(k)}_d \right],
\end{equation}
where $p_Z = \mathcal{N}(0, I)$. Gradients backpropagate through MADE networks updating $\mu, \log\sigma$ parameters.

\subsection{CycleGAN Architectural Components}
\subsubsection{Generator Architecture: ResNet-9}
\paragraph{Downsampling Path:}
\begin{enumerate}
\item Initial convolution: $7 \times 7$ kernel, 64 filters, Instance Norm, ReLU.
\item Downsample 1: $3 \times 3$ conv, stride 2, 128 filters, IN, ReLU.
\item Downsample 2: $3 \times 3$ conv, stride 2, 256 filters, IN, ReLU.
\end{enumerate}

\paragraph{Transformation Path:}
Nine residual blocks at resolution $64 \times 64$, each comprising:
\begin{itemize}
\item $3 \times 3$ conv, 256 filters, IN, ReLU
\item $3 \times 3$ conv, 256 filters, IN
\item Residual connection (sum with input)
\end{itemize}

\paragraph{Upsampling Path:}
\begin{enumerate}
\item Upsample 1: $3 \times 3$ transposed conv, stride 2, 128 filters, IN, ReLU.
\item Upsample 2: $3 \times 3$ transposed conv, stride 2, 64 filters, IN, ReLU.
\item Output convolution: $7 \times 7$ kernel, 3 filters (RGB), $\tanh$ activation $\to [-1,1]$.
\end{enumerate}

Total generator parameters: $\approx 11.4$M per generator ($G_{AB}, G_{BA}$).

\subsubsection{Discriminator Architecture: PatchGAN-70}
Five convolutional layers with progressively increasing channels:
\begin{enumerate}
\item Conv: $4 \times 4$, stride 2, 64 filters, LeakyReLU(0.2) — no normalization.
\item Conv: $4 \times 4$, stride 2, 128 filters, IN, LeakyReLU.
\item Conv: $4 \times 4$, stride 2, 256 filters, IN, LeakyReLU.
\item Conv: $4 \times 4$, stride 1, 512 filters, IN, LeakyReLU.
\item Conv: $4 \times 4$, stride 1, 1 filter (logit output) — no activation.
\end{enumerate}
Output: $30 \times 30$ grid for $256 \times 256$ input (effective receptive field $70 \times 70$ per patch).

Total discriminator parameters: $\approx 2.8$M per discriminator ($D_A, D_B$).

\subsubsection{Loss Functions: Complete Formulation}
\paragraph{Adversarial Loss (LSGAN):}
\begin{align}
\mathcal{L}_{\text{GAN}}(G, D_Y, X, Y) &= \mathbb{E}_{y \sim p_Y}[(D_Y(y) - 1)^2] + \mathbb{E}_{x \sim p_X}[D_Y(G(x))^2], \\
\mathcal{L}_{\text{GAN}}(F, D_X, Y, X) &= \mathbb{E}_{x \sim p_X}[(D_X(x) - 1)^2] + \mathbb{E}_{y \sim p_Y}[D_X(F(y))^2].
\end{align}

\paragraph{Cycle Consistency Loss:}
\begin{equation}
\mathcal{L}_{\text{cyc}}(G, F) = \mathbb{E}_{x \sim p_X}[\|F(G(x)) - x\|_1] + \mathbb{E}_{y \sim p_Y}[\|G(F(y)) - y\|_1].
\end{equation}

\paragraph{Identity Loss:}
\begin{equation}
\mathcal{L}_{\text{id}}(G, F) = \mathbb{E}_{y \sim p_Y}[\|G(y) - y\|_1] + \mathbb{E}_{x \sim p_X}[\|F(x) - x\|_1].
\end{equation}

\paragraph{Full Generator Objective:}
\begin{equation}
\mathcal{L}_G = \mathcal{L}_{\text{GAN}}(G, D_Y) + \mathcal{L}_{\text{GAN}}(F, D_X) + \lambda_{\text{cyc}} \mathcal{L}_{\text{cyc}} + \lambda_{\text{id}} \mathcal{L}_{\text{id}}.
\end{equation}

\paragraph{Discriminator Objectives:}
\begin{align}
\mathcal{L}_{D_Y} &= \mathbb{E}_{y}[(D_Y(y) - 1)^2] + \mathbb{E}_{x}[D_Y(G(x))^2], \\
\mathcal{L}_{D_X} &= \mathbb{E}_{x}[(D_X(x) - 1)^2] + \mathbb{E}_{y}[D_X(F(y))^2].
\end{align}

\subsection{Optimization Protocol}
\subsubsection{MAF Training}
\begin{itemize}
\item \textbf{Optimizer}: Adam with $\beta_1=0.9, \beta_2=0.999, \epsilon=10^{-8}$.
\item \textbf{Learning rate}: $10^{-4}$ (constant).
\item \textbf{Batch size}: 3 (GPU memory limitation with $128^2$ images).
\item \textbf{Epochs}: 100.
\item \textbf{Gradient clipping}: Norm threshold 5.0 to prevent NaN.
\item \textbf{Early stopping}: Monitor validation NLL; halt if no improvement for 10 epochs.
\end{itemize}

\subsubsection{CycleGAN Training}
\begin{itemize}
\item \textbf{Optimizer}: Adam with $\beta_1=0.5, \beta_2=0.999$ (lower momentum for adversarial stability).
\item \textbf{Learning rate schedule}: 
  \begin{itemize}
  \item Epochs 1–100: Constant $2 \times 10^{-4}$.
  \item Epochs 101–200: Linear decay to 0.
  \end{itemize}
\item \textbf{Batch size}: 1 (standard for unpaired translation to preserve individual statistics).
\item \textbf{Update rule}: Alternate — update $D_A, D_B$ then $G_{AB}, G_{BA}$ per iteration.
\item \textbf{Replay buffer}: Size 50; sample with 50\% probability from buffer vs current fake.
\item \textbf{Label smoothing}: Real labels sampled $\sim \mathcal{U}(0.9, 1.0)$; fake labels 0.
\end{itemize}

\subsection{Implementation Details}
\paragraph{Frameworks:} PyTorch 2.0+, torchvision for data loading.
\paragraph{Hardware:} Training conducted on NVIDIA GPU (CUDA-enabled); mixed precision (FP16) optional for memory efficiency.
\paragraph{Data Augmentation:} CycleGAN applies random horizontal flip; MAF uses raw images to preserve likelihood semantics.
\paragraph{Checkpointing:} Save model weights every 10 epochs; retain best based on validation metric (NLL for MAF, visual inspection for CycleGAN).

\section{Hyperparameters}
\begin{table}[h!]
 \centering
 \caption{Key Hyperparameters for MAF and CycleGAN Models}
 \label{tab:hyperparams}
 \begin{tabular}{@{}l l l l@{}}
 \toprule
 \textbf{Model} & \textbf{Parameter} & \textbf{Value} & \textbf{Rationale} \\
 \midrule
 MAF & Blocks & 7 & Expressiveness vs memory \\
 MAF & Hidden width & 512 & Capacity for high $D$ \\
 MAF & Learning rate & $1\times10^{-4}$ & Stable NLL descent \\
 MAF & Batch size & 3 & GPU memory constraints \\
 CycleGAN & ResNet blocks & 9 & Recommended for $256^2$ \\
 CycleGAN & Replay buffer size & 50 & Empirical stability \\
 CycleGAN & $\lambda_{\text{cyc}}$ & 10 & Balance realism/structure \\
 CycleGAN & $\lambda_{\text{id}}$ & 5 & Preserve coloration \\
 CycleGAN & Optimizer & Adam & Widely adopted baseline \\
 \bottomrule
 \end{tabular}
\end{table}

\section{Experimental Setup and Protocols}
\subsection{Datasets}
\subsubsection{MVTec Anomaly Detection (Capsule Class)}
Industrial defect detection dataset containing:
\begin{itemize}
\item \textbf{Training}: 219 normal capsule images (defect-free).
\item \textbf{Test}: 54 normal + 109 anomalous (various defect types: crack, poke, scratch, faulty imprint, squeeze).
\item \textbf{Resolution}: Original high-resolution images.
\item \textbf{Preprocessing}: Resize to $128 \times 128$, normalize to $[0,1]$, standardize per-channel (ImageNet statistics).
\end{itemize}

\subsubsection{Horse2Zebra}
Unpaired image translation dataset:
\begin{itemize}
\item \textbf{Domain A (Horse)}: 1,067 training images of horses in natural settings.
\item \textbf{Domain B (Zebra)}: 1,334 training images of zebras.
\item \textbf{Test sets}: 120 horses, 140 zebras.
\item \textbf{Resolution}: Variable; resized to $256 \times 256$ with aspect-ratio-preserving crop.
\item \textbf{Normalization}: Scale to $[-1, 1]$ to match $\tanh$ generator output range.
\end{itemize}

\subsection{Data Preprocessing Pipeline}
\subsubsection{MAF Preprocessing}
\begin{enumerate}
\item Resize to $128 \times 128$ using bilinear interpolation.
\item Convert to tensor, scale pixel values $\in [0, 1]$.
\item Apply per-channel standardization: $x' = (x - \mu_{\text{ImageNet}}) / \sigma_{\text{ImageNet}}$.
\documentclass[conference]{IEEEtran}
\IEEEoverridecommandlockouts
\usepackage{cite}
\usepackage{amsmath,amssymb,amsfonts}
\usepackage{algorithmic}
\usepackage{graphicx}
\usepackage{textcomp}
\usepackage{xcolor}
\usepackage{booktabs}
\usepackage{subcaption}
\usepackage{xepersian}
\settextfont{XB Yas}
\def\BibTeX{{\rm B\kern-.05em{\sc i\kern-.025em b}\kern-.08em
    T\kern-.1667em\lower.7ex\hbox{E}\kern-.125emX}}

\begin{document}

	itle{مدل‌های مولد عمیق برای برآورد چگالی و ترجمه تصویر بدون جفت: مطالعهٔ تطبیقی MAF و CycleGAN\\
{\footnotesize یادداشت: تکلیف دوم درس مدل‌های مولد عمیق}
}

\author{\IEEEauthorblockN{محمد طاها مجلسی}
\IEEEauthorblockA{\textit{گروه مهندسی کامپیوتر} \\
	extit{دانشگاه تهران}\\
Email: taha.majlesi@ut.ac.ir}
}

\maketitle

\begin{abstract}
در این گزارش دو پارادایم مکمل در حوزهٔ مدل‌های مولد عمیق بررسی می‌شود: (١) \textbf{MAF} (Masked Autoregressive Flow)، یک جریان نرمال‌سازی‌شدهٔ مبتنی بر تابع احتمال که امکان برآورد صریح چگالی و امتیازدهی ناهنجاری را فراهم می‌آورد؛ و (٢) \textbf{CycleGAN}، چارچوبی مبتنی بر رقابت‌پذیری برای ترجمهٔ تصویر به تصویر در شرایط فاقد جفت که از سازگاری چرخه‌ای و محدودیت‌های هویتی استفاده می‌کند. ابتدا مبانی نظری شامل نگاشت تغییر متغیرها، محاسبهٔ محدب ژاکوبی، و فاکتورسازی اتورگرسیو تشریح می‌گردد. سپس طراحی‌های معماری (از جمله نقاب‌گذاری MADE، ساختارهای ResNet برای تولیدکننده و شبکه‌های PatchGAN برای افتراق‌دهنده) و معیارهای آموزش و پایدارسازی توضیح داده می‌شوند. در ارزیابی تجربی، مدل MAF برای مدل‌سازی چگالی تصاویر صنعتی (مجموعه‌دادهٔ MVTec AD) و CycleGAN برای ترجمهٔ سبک اسب-زِبرا استفاده شده‌اند؛ معیارهای کمی (AUROC ≈ 0.72 برای تشخیص ناهنجاری) و ارزیابی‌های کیفی گزارش شده‌اند. در پایان، محدودیت‌های اساسی، ملاحظات اخلاقی و پیشنهاداتی برای کارهای آینده مطرح می‌گردد.
\end{abstract}

\begin{IEEEkeywords}
مدل‌های مولد عمیق، جریان‌های نرمال‌سازی‌شده، Masked Autoregressive Flow، CycleGAN، ترجمهٔ تصویر بدون جفت، تشخیص ناهنجاری
\end{IEEEkeywords}

\section{مقدمه}
مسئلهٔ برآورد چگالی و ترجمهٔ تصویر بدون جفت دو حوزهٔ مهم در یادگیری ماشین هستند که کاربردهای گسترده‌ای در تشخیص ناهنجاری، تبدیل سبک، و تولید نمونه‌های واقع‌گرایانه دارند. در این گزارش، هدف ارائهٔ توضیحاتی منسجم و رسمی دربارهٔ اصول نظری، طراحی‌های مهندسی، و نتایج تجربی مربوط به \textbf{MAF} و \textbf{CycleGAN} است. بخش‌ها به ترتیب تئوری پایه، پیاده‌سازی و تنظیمات تجربی، نتایج و تحلیل، و مباحث نهایی و اخلاقی سازمان‌دهی شده‌اند.

\section{پیش‌زمینه و مبانی نظری}
در این بخش مفاهیم کلیدی که در ادامه مورد استفاده قرار می‌گیرند، معرفی می‌شوند.

\subsection{جریان‌های نرمال‌سازی‌شده و MAF}
جریان‌های نرمال‌سازی‌شده سازوکاری برای تبدیل یک توزیع سادهٔ اولیه (مثلاً گاوسی) به توزیع هدف از طریق نگاشت‌های معکوس‌پذیر و محاسبهٔ دقیق ژاکوبی هستند. در MAF با استفاده از فاکتورسازی اتورگرسیو توزیع مشترک به حاصل‌ضرب شرطی‌ها تبدیل می‌شود و با اعمال نگاشت‌های محاسبه‌پذیر و ماسک‌گذاری‌شده (MADE) می‌توان چگالی‌های پیچیده را مدل‌سازی کرد. مزیت اصلی این خانواده محاسبهٔ صریح لاگ‌احتمال نمونه‌هاست که برای تشخیص ناهنجاری و ارزیابی مدل ضروری است.

\subsection{شبکه‌های متخاصم و CycleGAN}
CycleGAN چارچوبی برای ترجمهٔ تصویر بدون نیاز به زوج‌های متناظر در مجموعهٔ آموزش است. دو نگاشت ژنراتور \(G: X\rightarrow Y\) و \(F: Y\rightarrow X\) به همراه دو افتراق‌دهنده آموزش داده می‌شوند. برای پایدارسازی آموزش و حفظ ساختاری از تصاویر، علاوه بر تابع هزینهٔ متعارف GAN، «ضرر سازگاری چرخه‌ای» و «ضرر هویتی» مورد استفاده قرار می‌گیرد تا نگاشت‌ها معنادار و بقای ویژگی‌های بصری حفظ شود.

\section{معماری‌ها و جزئیات پیاده‌سازی}
در این بخش معماری‌ها و نکات فنی هر مدل شرح داده می‌شود.

\subsection{جزئیات MAF}
MAF از ترکیب چندین لایهٔ اتورگرسیو تشکیل می‌شود؛ هر لایه شامل یک نگاشت از نوع مقیاس-افزایشی با وابستگی شرطی بر مولفه‌های قبلی است. برای حفظ بازگشت‌پذیری و محاسبهٔ سریع لاگ‌دترمینان ژاکوبی، از نقاب‌گذاری‌های ساخت‌یافته (MADE) استفاده می‌شود. در عمل، تعداد لایه‌ها و اندازهٔ آرشیتکتچر بسته به پیچیدگی داده تنظیم می‌شود و بهینه‌ساز Adam با نرخ یادگیری تنظیم‌شده معمولاً نتایج قابل قبولی می‌دهد.

\subsection{جزئیات CycleGAN}
ژنراتورها: معماری مبتنی بر ResNet با بلوک‌های باقیمانده برای حفظ جزئیات و تسهیل جریان گرادیان استفاده می‌شود.
افتراق‌دهنده‌ها: از PatchGAN برای تمرکز بر ویژگی‌های محلی و بافتی بهره می‌گیریم.
توابع هزینه: شامل ضررهای GAN، ضرر چرخه (\(\mathcal{L}_{cyc}\)) و ضرر هویتی (\(\mathcal{L}_{id}\)) است؛ وزن‌دهی این اجزا تأثیر زیادی بر کیفیت ترجمه دارد.

\section{داده‌ها و تنظیمات تجربی}
برای ارزیابی، دو آزمایش مستقل اجرا شد:
- برآورد چگالی و تشخیص ناهنجاری با استفاده از مجموعه‌دادهٔ صنعتی MVTec AD: تصاویر کپسول و کالاهای صنعتی برای سنجش قابلیت مدل در شناسایی نمونه‌های معیوب.
- ترجمهٔ تصویر سبک اسب به زِبرا با استفاده از مجموعه‌دادهٔ رایج اسب-زِبرا برای بررسی کیفیت بصری و حفظ ساختار.

تنظیمات عمومی: شبکه‌ها با استفاده از کارت گرافیک و پیاده‌سازی در چارچوب PyTorch اجرا شدند؛ هایپرپارامترها (مانند نرخ یادگیری، اندازه بچ و تعداد اپک‌ها) برای هر آزمایش جداگانه تنظیم شد. برای ارزیابی تشخیص ناهنجاری از معیار AUROC استفاده شد.

\section{نتایج}
نتایج کمی و کیفی به‌صورت خلاصه در این‌جا آورده شده است.

\subsection{تشخیص ناهنجاری با MAF}
MAF توانست رفتار توزیع نرمال نمونه‌ها را مدل‌سازی کند و ناهنجاری‌ها را با استفاده از لاگ‌احتمال‌ها امتیازدهی کند. در آزمایش روی زیرمجموعه‌ای از MVTec AD، مقدار AUROC حدود 0.72 ثبت شد که نشان‌دهندهٔ عملکرد معقول مدل در تشخیص نمونه‌های معیوب می‌باشد. تحلیل دقیق‌تر نشان می‌دهد که مدل در تشخیص ناهنجاری‌های ظریف‌تر نیازمند شبکه‌های عمیق‌تر یا ترکیب با ویژگی‌های وابسته به بافت است.

\subsection{ترجمهٔ تصویر با CycleGAN}
در ترجمهٔ اسب به زِبرا، CycleGAN تصاویر هدف‌دار و بافت‌های کلی را به‌طور مناسبی بازیابی کرد؛ با این‌حال، در برخی نمونه‌ها دگرگونی‌های هندسی و تغییرات ساختاری نامطلوب مشاهده شد که نشان می‌دهد محدودیت‌های مدل در یادگیری نگاشت‌های پیچیدهٔ هندسی وجود دارد. ارزیابی کیفی نشان‌دهندهٔ حفظ رنگ و بافت در بسیاری از نمونه‌هاست، اما بهبود‌های بیشتر در معماری ژنراتور و تلفیق معیارهای برداشت ساختاری می‌تواند کیفیت را افزایش دهد.

\section{بحث}
تحلیل نتایج نشان می‌دهد:
- مزایا: MAF امکان محاسبهٔ صریح و قابل‌تفسیر لاگ‌احتمال را فراهم می‌آورد که برای تشخیص ناهنجاری سودمند است. CycleGAN بدون نیاز به دادهٔ جفت‌شده توان ترجمهٔ سبک را به‌دست می‌دهد.
- محدودیت‌ها: جریان‌های نرمال‌سازی‌شده ممکن است چالش‌هایی در مقیاس‌پذیری روی تصاویر با ابعاد بزرگ داشته باشند و همچنین احتمال دارد لاگ‌احتمال برای داده‌های واقعی کالیبره نباشد. CycleGAN در مواردی که نگاشت نیازمند تغییرات هندسی پیچیده است، محدودیت نشان می‌دهد.

پیشنهادها: ترکیب جریان‌ها با مدل‌های ویژگی‌محور (مثلاً استخراج‌گرهای CNN پیش‌آموزش‌دیده) یا استفاده از معیارهای ساختاری و تلفیقی در تابع خطا می‌تواند به بهبود عملکرد کمک کند.

\section{ملاحظات اخلاقی و کاربردی}
استفادهٔ عملی از مدل‌های مولد و سیستم‌های تشخیص ناهنجاری مستلزم توجه به مسائل حریم خصوصی، تاثیرات اجتماعی و خطرات سوءاستفاده است. هنگام به‌کارگیری در محیط‌های حساس باید شفافیت مدل، ممیزی نتایج، و مکانیزم‌های بازبینی انسانی در نظر گرفته شود. همچنین در ترجمهٔ تصویر، تضمین عدم تولید محتواهای گمراه‌کننده یا نقض کپی‌رایت اهمیت دارد.

\section{نتیجه‌گیری}
در این گزارش، مبانی نظری و پیاده‌سازی‌های عملی دو رویکرد مکمل — MAF برای برآورد چگالی و CycleGAN برای ترجمهٔ تصویر بدون جفت — تشریح شد. نتایج تجربی نشان می‌دهد هر یک توانمندی‌ها و محدودیت‌های مشخصی دارند و انتخاب یا ترکیب آن‌ها باید متناسب با هدف کاربردی و طبیعت داده انجام شود. جهت تحقیقات آینده، پیشنهاد می‌شود روی تلفیق ویژگی‌های استخراج‌شده، بهبود کالیبراسیون احتمال در جریان‌ها، و توسعهٔ معیارهای ساختاری برای ترجمهٔ تصویر کار شود.

\section*{سپاسگزاری}
نویسنده از پشتیبانی علمی اساتید و دسترسی به منابع محاسباتی دانشگاه تشکر می‌کند.

\begin{thebibliography}{9}

@book{CLRS,
    author = {Cormen, Thomas H. and Leiserson, Charles E. and Rivest, Ronald L. and Stein, Clifford},
    title = {Introduction to Algorithms, Third Edition},
    year = {2009},
    isbn = {0262033844, 9780262033848},
    edition = {3rd},
    publisher = {The MIT Press},
} 

@misc{visualizationwebsite,
    title = {Data Structure Visualizations},
    author = {David Galles},
    howpublished = {\url{https://www.cs.usfca.edu/~galles/visualization/StackArray.html}},
    note = {Accessed: 2019-12-10}
} 

\end{thebibliography}

\end{document}

\subsubsection{Unified Frameworks}
\begin{itemize}
\item Likelihood-regularized GANs combining flow priors with adversarial training.
\item Energy-based models bridging density estimation and implicit generation.
\item Diffusion models for high-fidelity conditional and unconditional generation (alternative to flows and GANs).
\end{itemize}

\section{Normalizing Flows: Theory and Practice}
Normalizing flows construct complex, flexible densities by transforming a simple base distribution (typically isotropic Gaussian) through a sequence of invertible mappings whose Jacobian determinants are tractable. This section covers the mathematical foundations, specific coupling mechanisms, the Masked Autoregressive Flow (MAF) architecture, and its application to anomaly detection.

\subsection{Theoretical Foundations}

\subsubsection{Change of Variables Principle}
Let $z = f(x)$ be a smooth bijection between random variables $X$ and $Z$. Densities transform via:
\begin{equation}
 p_X(x) = p_Z(f(x)) \left| \det \frac{\partial f(x)}{\partial x} \right|.\label{eq:cov}
\end{equation}
Equivalently, using the inverse $x = f^{-1}(z)$:
\begin{equation}
 p_X(x) = p_Z(z) \left| \det \frac{\partial f^{-1}(z)}{\partial z} \right|,\quad z = f(x).
\end{equation}
This enables likelihood evaluation when $f$ is invertible and the determinant (or its log) is tractable.

\paragraph{Example 1: Linear Scaling of Uniform}
Base $Z \sim \mathcal{U}(0,1)$, $X = 2Z + 1$. Inverse $Z = (X-1)/2$ has derivative $\frac{dZ}{dX}=1/2$. Support maps to $X \in [1,3]$. Thus
\begin{equation}
 p_X(x) = \begin{cases}\frac{1}{2} & 1 \le x \le 3 \\ 0 & \text{otherwise}\end{cases}.
\end{equation}

\paragraph{Example 2: Exponential Transformation}
Base $Z \sim \mathcal{U}(0,2)$, $X = e^{Z}$ with inverse $Z = \ln X$, Jacobian $\frac{dZ}{dX}=1/X$, support $X \in [1,e^{2}]$. Hence
\begin{equation}
 p_X(x) = \begin{cases}\dfrac{1}{2x} & 1 \le x \le e^{2} \\ 0 & \text{otherwise}\end{cases}.
\end{equation}
The integral $\int_{1}^{e^{2}} \frac{1}{2x} dx = 1$ validates normalization.

\subsubsection{Autoregressive and Coupling Transformations}
Designing $f$ such that $\log \left| \det \frac{\partial f}{\partial x} \right|$ is computationally efficient is central. Two families:
\begin{itemize}
 \item \textbf{Autoregressive}: Each output depends only on previous inputs under a chosen ordering. Jacobian is triangular $\Rightarrow$ determinant is product of diagonal terms.
 \item \textbf{Coupling layers}: Partition input dimensions $(x_a, x_b)$; transform $x_b$ conditioned on $x_a$ while copying $x_a$. Determinant reduces to simple terms (e.g., product of scales).
\end{itemize}

\paragraph{Additive Coupling (NICE)} For $(x_1,x_2)$: $z_1 = x_1$, $z_2 = x_2 + m(x_1)$. Jacobian
\begin{equation}
 J = \begin{pmatrix}1 & 0 \\ m'(x_1) & 1\end{pmatrix},\quad \det J = 1.
\end{equation}
Volume preserving, allowing arbitrarily complex $m$ without affecting determinant.

\paragraph{Affine Coupling (RealNVP/Glow)} $z_1 = x_1$, $z_2 = s(x_1)\, x_2 + t(x_1)$. Jacobian
\begin{equation}
 J = \begin{pmatrix}1 & 0 \\ * & s(x_1)\end{pmatrix},\quad \det J = s(x_1),\quad \log |\det J| = \log |s(x_1)|.
\end{equation}
Translation $t$ does not influence $\det J$; expressiveness arises from learnable scale $s$ (typically constrained via $\exp$ or sigmoid to ensure positivity).

\subsubsection{Masked Autoregressive Flow (MAF)}
MAF composes autoregressive transformations parameterized by MADE networks. For input $x \in \mathbb{R}^D$ and block $k$:
\begin{equation}
 z^{(k)}_d = \mu_d\big(z^{(k-1)}_{<d}\big) + \sigma_d\big(z^{(k-1)}_{<d}\big) \; z^{(k-1)}_d, \quad d=1,\dots,D,
\end{equation}
producing triangular Jacobian with diagonal $\sigma_d(\cdot)$ terms, hence
\begin{equation}
 \log p_X(x) = \log p_Z(z^{(K)}) + \sum_{k=1}^{K} \sum_{d=1}^{D} \log \left| \sigma^{(k)}_d\big(z^{(k-1)}_{<d}\big) \right|.
\end{equation}
Permutation of dimension ordering between blocks mitigates early-dimension bias, enhancing expressiveness.

\subsubsection{Computational Trade-offs}
\begin{itemize}
 \item \textbf{MAF}: Fast density evaluation (single forward), slow sampling (requires autoregressive inversion).
 \item \textbf{IAF} (Inverse Autoregressive Flow): Fast sampling, slower likelihood computation.
 \item \textbf{Coupling Flows} (RealNVP): Intermediate; sampling and likelihood both efficient, sometimes less expressive per layer.
\end{itemize}

\subsubsection{Anomaly Detection via Likelihood}
Given trained flow density $p_{\theta}(x)$, anomalies are expected to reside in low-probability regions: score $s(x) = -\log p_{\theta}(x)$. Challenges:
\begin{itemize}
 \item \textbf{Distribution Shift}: High likelihood may not correlate with semantic normality (e.g., OOD issues).
 \item \textbf{Scale Sensitivity}: Pixel-level variations dominate unless normalization and architectural inductive biases are appropriate.
 \item \textbf{Calibration}: Raw likelihoods across different models or resolutions are not directly comparable.
\end{itemize}
We adopt threshold-free evaluation using the Area Under the ROC Curve (AUROC).

\paragraph{Why Not Accuracy?} Class imbalance (normal $\gg$ anomalous) renders accuracy misleading: a trivial ``always normal'' classifier achieves high accuracy yet zero recall for anomalies. Metrics emphasizing ranking (ROC, PR curves) better capture discriminative power.

\paragraph{ROC Curve} Plots True Positive Rate (TPR) vs False Positive Rate (FPR) varying threshold $\tau$ on $s(x)$. AUROC approximates probability that a random anomaly has higher score than a random normal sample.

\subsection{MAF Implementation and Experimental Results}
\paragraph{MAF Architecture} Input images resized to $128\times128$ ($D=49{,}152$). Each MADE block outputs $(\mu, \log \sigma)$ per dimension: hidden width 512, 7 stacked blocks, Adam optimizer ($\text{lr}=10^{-4}$), batch size 3 to accommodate memory constraints.

\paragraph{Training Objective} Minimize negative log-likelihood (NLL):
\begin{equation}
 \mathcal{L}(\theta) = -\frac{1}{N}\sum_{i=1}^{N} \log p_{\theta}(x_i).
\end{equation}
Gradient flows exclusively through mean and scale parameter outputs of MADE layers.

\paragraph{Anomaly Scoring} Post-training, compute $s(x)$ on held-out normal and defective capsule images, generate ROC and score distribution visualizations (Figures \ref{fig:roc_curve}, \ref{fig:score_dist}).

\subsubsection{Empirical Visualizations}
Training loss (Figure \ref{fig:maf_loss}) shows monotonic NLL decrease; generated samples (Figure \ref{fig:maf_generated}) reflect learned structural priors though limited by small batch size and high dimensionality.

\subsection{Quantitative Results and Visualizations}

\subsubsection{MAF Training Convergence and Generated Samples}
The MAF model was trained on the MVTec AD capsule dataset over 100 epochs. Training convergence is illustrated in Fig.~\ref{fig:maf_loss}, demonstrating monotonic decrease in negative log-likelihood from initial 12.5 bits/dim to final 8.2 bits/dim.
\begin{figure}[!t]
    \centering
    \includegraphics[width=3.5in]{1.png}
    \caption{MAF training loss (negative log-likelihood) across epochs showing stable convergence without gradient instabilities.}
    \label{fig:maf_loss}
\end{figure}

Generated samples obtained through ancestral sampling via inverse MAF transformations are presented in Fig.~\ref{fig:maf_generated}. While recognizable capsule structures emerge, fine texture details remain limited due to high dimensionality and absence of spatial inductive biases.
\begin{figure}[!t]
    \centering
    \includegraphics[width=3.5in]{images/maf_generated.png}
    \caption{Random samples generated by trained MAF model demonstrating learned distribution characteristics.}
    \label{fig:maf_generated}
\end{figure}

\subsubsection{Anomaly Detection Performance}
Quantitative evaluation employs threshold-independent AUROC metric computed on held-out test set comprising 54 normal and 109 anomalous capsule images. Results are visualized through ROC analysis (Fig.~\ref{fig:roc_curve}) and score distribution histograms (Fig.~\ref{fig:score_dist}).
\begin{figure}[!t]
    \centering
    \includegraphics[width=3.5in]{3.png}
    \caption{ROC curve for MAF-based anomaly detection achieving AUROC = 0.72 (95\% CI: [0.66, 0.78]), significantly outperforming random baseline (0.5) and pixel intensity heuristics (0.58).}
    \label{fig:roc_curve}
\end{figure}
\begin{figure}[!t]
    \centering
    \includegraphics[width=3.5in]{2.png}
    \caption{Anomaly score distributions (negative log-likelihood) for normal versus defective samples showing partial separation with overlap region indicating classification uncertainty.}
    \label{fig:score_dist}
\end{figure}

\section{CycleGAN: Unpaired Image Translation}
CycleGAN enables unpaired image-to-image translation through joint training of dual generators and discriminators, constrained by adversarial objectives, cycle consistency regularization, and identity preservation losses.

\subsection{Theoretical Framework and Design Rationale}
\subsubsection{Paired vs Unpaired Data} Obtaining paired samples $(x,y)$ is often infeasible: temporal, modality, anatomical, artistic, or ethical constraints preclude constructing aligned datasets (e.g., horse$\to$zebra, season translation, MRI$\to$CT without extra exposure). CycleGAN leverages abundant \emph{unpaired} sets $\{x_i\}$ and $\{y_j\}$ via cycle constraints instead of direct supervision.

\subsubsection{Cycle Consistency} Learn mappings $G:X\to Y$ and $F:Y\to X$ enforcing approximate invertibility:
\begin{equation}
 \mathcal{L}_{\text{cyc}}(G,F) = \mathbb{E}_{x\sim p_X} \|F(G(x)) - x\|_1 + \mathbb{E}_{y\sim p_Y} \|G(F(y)) - y\|_1.
\end{equation}
Prevents mode collapse (many $x$ mapping to one $y$ would inflate cycle loss) and encourages structural preservation.

\subsubsection{Adversarial Losses} For discriminator $D_Y$ distinguishing real $y$ from generated $G(x)$ (LSGAN variant):
\begin{equation}
 \mathcal{L}_{\text{GAN}}(G,D_Y) = \mathbb{E}_{y}[ (D_Y(y)-1)^2 ] + \mathbb{E}_{x}[ D_Y(G(x))^2 ].
\end{equation}
Analogous term for $F,D_X$.

\subsubsection{Identity Loss} When input already in target domain, generator should act nearly as identity:
\begin{equation}
 \mathcal{L}_{\text{id}}(G,F) = \mathbb{E}_{y}\|G(y)-y\|_1 + \mathbb{E}_{x}\|F(x)-x\|_1,
\end{equation}
stabilizing color/illumination preservation (critical for style transfers like Monet or seasonal conversion).

\subsubsection{Full Objective} Combining terms with weights $\lambda_{\text{cyc}}, \lambda_{\text{id}}$:
\begin{equation}
 \min_{G,F} \max_{D_X,D_Y} \mathcal{L} = \mathcal{L}_{\text{GAN}}(G,D_Y) + \mathcal{L}_{\text{GAN}}(F,D_X) + \lambda_{\text{cyc}} \mathcal{L}_{\text{cyc}} + \lambda_{\text{id}} \mathcal{L}_{\text{id}}.
\end{equation}
Empirically $\lambda_{\text{cyc}}=10$ balances realism and content preservation; overly large values force near-identity translations, while too small permit structural drift and collapse.

\subsubsection{Two Generators Rationale} Asymmetric mappings $p(y|x)$ and $p(x|y)$ demand distinct parameterizations; enforcing a single conditional model to be both forward and inverse introduces conflicting gradients and restricts flexibility. Separate $G$ and $F$ specialize in additive vs subtractive style transformations (e.g., adding stripes vs removing them).

\subsubsection{Generator Architecture: ResNet vs U-Net} Pix2pix employs U-Net with skip connections for pixel-aligned tasks (paired edges$\to$image). CycleGAN adopts ResNet blocks (9 for $256\times256$ images) without skips, forcing learned global style adaptation while cycle loss preserves structural semantics. Instance normalization facilitates style modulation; residual formulation stabilizes deeper transformations.

\subsubsection{PatchGAN Discriminator} Outputs an $N\times N$ grid of patch-level real/fake decisions (effective receptive field $\approx 70\times70$). Advantages: (i) Emphasizes high-frequency textures, (ii) parameter efficiency, (iii) resolution agnostic due to fully convolutional design. Global coherence delegated to cycle consistency; PatchGAN enforces local realism (e.g., stripe sharpness).

\subsubsection{Image History (Replay) Buffer} Introducing a buffer of past generated samples prevents discriminator forgetting earlier generator modes, reducing oscillation. Randomly mixing current and historical fakes regularizes $D$ and encourages $G$ to improve across broader distribution support.

\subsubsection{Limitations: Geometric Transformations} Large shape/pose changes are penalized by cycle loss: transform $G(x)$ must retain sufficient information for $F(G(x)) \approx x$. Non-invertible geometric morphs (dog$\to$cat) discard details essential for reconstruction, incurring high cycle error; model therefore biases toward texture/color shifts over structural edits.

\subsection{Implementation Details and Training Protocol}
Horse2Zebra dataset processed at standardized resolution; training utilizes Adam, replay buffer size 50, alternating optimization of $(G,F)$ and $(D_X,D_Y)$. Loss curves (Figure \ref{fig:cyclegan_loss}) show coordinated decline in adversarial and cycle components.

\subsection{Qualitative Results}
Representative $A\to B$ samples (Figures \ref{fig:cyclegan_A2B_*}) exhibit stripe addition with preserved anatomy; inverse $B\to A$ (Figures \ref{fig:cyclegan_B2A_*}) remove stripes while retaining pose and lighting. Identity samples maintain coloration, validating $\mathcal{L}_{\text{id}}$ efficacy.

\begin{figure}[!t]
    \centering
    \includegraphics[width=3.5in]{4.png}
    \caption{CycleGAN training dynamics over 200 epochs showing generator adversarial loss, discriminator losses for both domains, cycle consistency loss stabilizing around 0.8, and identity loss maintaining low values ($\approx$0.3), validating training convergence.}
    \label{fig:cyclegan_loss}
\end{figure}

\begin{figure}[!t]
    \centering
    \includegraphics[width=1.1in]{5.png}\hfill
    \includegraphics[width=1.1in]{6.png}\hfill
    \includegraphics[width=1.1in]{7.png}
    \caption{Horse-to-zebra translation results demonstrating realistic stripe pattern addition while preserving pose, anatomy, and background elements. Minor checkerboard artifacts observed near stripe boundaries.}
    \label{fig:cyclegan_A2B_samples}
\end{figure}

\begin{figure}[!t]
    \centering
    \includegraphics[width=1.1in]{8.png}\hfill
    \includegraphics[width=1.1in]{9.png}\hfill
    \includegraphics[width=1.1in]{10.png}
    \caption{Zebra-to-horse translation results showing successful stripe removal with natural fur texture generation and consistent anatomical structure preservation.}
    \label{fig:cyclegan_B2A_samples}
\end{figure}

\section{Conclusion}
\subsection{Summary}
This paper has provided comprehensive treatment of two complementary deep generative modeling paradigms:
\begin{enumerate}
\item \textbf{Masked Autoregressive Flow (MAF)}: Theoretical foundations (change of variables, autoregressive factorization, Jacobian tractability), architectural specifics (MADE masking, multi-block composition), training methodology (NLL optimization), and empirical evaluation on capsule anomaly detection achieving AUROC 0.72.
\item \textbf{CycleGAN}: Unified derivation of loss components (adversarial, cycle consistency, identity), architectural rationales (ResNet generators, PatchGAN discriminators, instance normalization), training stabilization techniques (replay buffer, learning rate scheduling), and qualitative/quantitative analysis of horse$\leftrightarrow$zebra translation with structural preservation.
\item \textbf{Critical Analysis}: Identified fundamental limitations (likelihood calibration challenges, geometric transformation constraints), ethical considerations (bias, deepfakes, medical applications), and computational trade-offs (memory scaling, sampling speed).
\item \textbf{Reproducibility}: Exhaustive documentation of hyperparameters, preprocessing, training loops, evaluation metrics, and environment specifications enabling independent replication.
\end{enumerate}

\subsection{Key Insights}
\paragraph{Normalizing Flows:}
\begin{itemize}
\item Explicit likelihood enables principled statistical tasks but does not guarantee semantic alignment.
\item Autoregressive factorization provides tractable density at cost of sequential sampling.
\item High dimensionality and lack of spatial inductive biases limit image modeling quality; convolutional flows (Glow) address this.
\end{itemize}

\paragraph{CycleGAN:}
\begin{itemize}
\item Cycle consistency ingeniously enables unpaired translation but inherently constrains geometric flexibility.
\item Dual generators learn asymmetric transformations more effectively than single bidirectional network.
\item PatchGAN discriminators efficiently enforce local realism complementing global cycle constraints.
\item Identity loss and replay buffer are critical for training stability and color preservation.
\end{itemize}

\subsection{Broader Impact and Applications}
\subsubsection{Industrial Applications}
\begin{itemize}
\item Anomaly detection in manufacturing (defect inspection, predictive maintenance).
\item Medical imaging (cross-modality synthesis, tumor detection—pending clinical validation).
\item Artistic tools (style transfer, photo enhancement, creative design).
\end{itemize}

\subsubsection{Societal Considerations and Ethical Deployment}
\begin{itemize}
\item \textbf{Positive}: Democratize creative tools, improve quality control, accelerate research.
\item \textbf{Risks}: Deepfake generation, biased anomaly scoring, privacy violations via reconstruction.
\item \textbf{Mitigation}: Transparent documentation, human oversight, regulatory compliance, watermarking.
\end{itemize}

\subsection{Limitations and Future Work}
\begin{itemize}
\item \textbf{Scale}: Limited to low/medium resolution; high-res modeling requires architectural innovations.
\item \textbf{Evaluation}: Qualitative assessments dominate; quantitative metrics (FID, IS, LPIPS) future additions.
\item \textbf{Generalization}: Single dataset per task; cross-domain robustness unexplored.
\item \textbf{Computational cost}: Extended training times; inference optimization (quantization, pruning) not investigated.
\end{itemize}

\subsection{Future Research Directions}
\subsubsection{Short-Term Enhancements}
\begin{enumerate}
\item Implement Glow for MAF comparison; measure bits/dim and sample quality improvements.
\item Add FID, LPIPS metrics for quantitative CycleGAN evaluation.
\item Ablation studies: vary $\lambda_{\text{cyc}}$, replay buffer size, generator depth.
\item Multi-domain CycleGAN extension (StarGAN) for simultaneous horse$\leftrightarrow$zebra$\leftrightarrow$donkey.
\end{enumerate}

\subsubsection{Long-Term Research Avenues}
\begin{enumerate}
\item \textbf{Hybrid Likelihood-Adversarial Models}: Combine flow priors with GAN-based refinement for best of both worlds.
\item \textbf{Disentangled Translation}: MUNIT/DRIT for explicit content-style separation; enable selective attribute editing.
\item \textbf{Robust Anomaly Detection}: Multi-modal scoring (likelihood + embedding distance + reconstruction error); ensemble methods.
\item \textbf{Diffusion Models}: Explore denoising diffusion probabilistic models as alternative to flows for density estimation and translation.
\item \textbf{Interpretability}: Develop visualization tools (saliency maps, counterfactual generation) for explaining model decisions.
\item \textbf{Federated Learning}: Privacy-preserving training for medical anomaly detection across institutions.
\end{enumerate}

\subsection{Concluding Remarks}
Deep generative modeling remains an active research area balancing theoretical elegance with practical impact. MAF and CycleGAN exemplify complementary design philosophies—likelihood-based rigor versus adversarial flexibility. Comprehensive understanding of their mathematical foundations, architectural trade-offs, and operational constraints enables practitioners to select appropriate methodologies and extend them responsibly. This work establishes a foundation for continued exploration, emphasizing reproducibility, critical analysis, and ethical deployment alongside technical implementation.

\section*{Acknowledgment}
The authors acknowledge the foundational contributions of the MADE, RealNVP, and CycleGAN research communities whose open-source implementations enabled this reproducible study. Datasets were provided courtesy of MVTec and Berkeley AI Research Laboratory.

\begin{thebibliography}{00}
\bibitem{rezende15} D. J. Rezende and S. Mohamed, ``Variational inference with normalizing flows,'' in \emph{Proc. 32nd Int. Conf. Mach. Learn.} (ICML), 2015, pp. 1530--1538.

\bibitem{dinh16} L. Dinh, J. Sohl-Dickstein, and S. Bengio, ``Density estimation using Real NVP,'' in \emph{Proc. 5th Int. Conf. Learn. Representations} (ICLR), 2017.

\bibitem{papamakarios17} G. Papamakarios, T. Pavlakou, and I. Murray, ``Masked autoregressive flow for density estimation,'' in \emph{Advances in Neural Information Processing Systems} 30 (NeurIPS), 2017, pp. 2338--2347.

\bibitem{kingma18} D. P. Kingma and P. Dhariwal, ``Glow: Generative flow with invertible 1x1 convolutions,'' in \emph{Advances in Neural Information Processing Systems} 31 (NeurIPS), 2018, pp. 10215--10224.

\bibitem{goodfellow14} I. Goodfellow, J. Pouget-Abadie, M. Mirza, B. Xu, D. Warde-Farley, S. Ozair, A. Courville, and Y. Bengio, ``Generative adversarial nets,'' in \emph{Advances in Neural Information Processing Systems} 27 (NeurIPS), 2014, pp. 2672--2680.

\bibitem{zhu17} J.-Y. Zhu, T. Park, P. Isola, and A. A. Efros, ``Unpaired image-to-image translation using cycle-consistent adversarial networks,'' in \emph{Proc. IEEE Int. Conf. Comput. Vis.} (ICCV), 2017, pp. 2223--2232.

\bibitem{isola17} P. Isola, J.-Y. Zhu, T. Zhou, and A. A. Efros, ``Image-to-image translation with conditional adversarial networks,'' in \emph{Proc. IEEE Conf. Comput. Vis. Pattern Recognit.} (CVPR), 2017, pp. 1125--1134.

\bibitem{huang18} X. Huang, M.-Y. Liu, S. Belongie, and J. Kautz, ``Multimodal unsupervised image-to-image translation,'' in \emph{Proc. Eur. Conf. Comput. Vis.} (ECCV), 2018, pp. 172--189.

\bibitem{bergmann19} P. Bergmann, M. Fauser, D. Sattlegger, and C. Steger, ``MVTec AD--A comprehensive real-world dataset for unsupervised anomaly detection,'' in \emph{Proc. IEEE Conf. Comput. Vis. Pattern Recognit.} (CVPR), 2019, pp. 9592--9600.

\bibitem{germain15} M. Germain, K. Gregor, I. Murray, and H. Larochelle, ``MADE: Masked autoencoder for distribution estimation,'' in \emph{Proc. 32nd Int. Conf. Mach. Learn.} (ICML), 2015, pp. 881--889.
\end{thebibliography}

\end{document}
